\section{The R vocabulary}

\begin{frame}{A compact map of the functions}
\small
\begin{tabularx}{\textwidth}{p{2.3cm} X X}
\toprule
Task & Core functions & Question answered \\
\midrule
Start & \texttt{library()}, \texttt{read\_csv()} & What object did we load? \\
Inspect & \texttt{glimpse()}, \texttt{names()}, \texttt{count()} & What does a row represent? \\
Choose rows/columns & \texttt{filter()}, \texttt{select()} & Which observations and variables belong? \\
Construct & \texttt{mutate()}, \texttt{arrange()} & What new measure or ordering is useful? \\
Aggregate & \texttt{group\_by()}, \texttt{summarise()} & What statistic belongs to each group? \\
Visualize & \texttt{ggplot()}, \texttt{aes()}, \texttt{geom\_*()}, \texttt{labs()} & What comparison should the reader see? \\
\bottomrule
\end{tabularx}
\end{frame}

\begin{frame}[fragile]{Install once; attach each session}
\begin{lstlisting}[language=R]
# Run only if tidyverse is absent.
install.packages("tidyverse")

# Run at the start of each session that uses it.
library(tidyverse)
\end{lstlisting}
\vspace{0.35cm}
Installation puts a package on the computer. \texttt{library()} makes its functions available in the current R session.
\end{frame}

\begin{frame}[fragile]{The file format determines the reader}
\begin{lstlisting}[language=R]
health <- read_csv("data/lesson-4/health-2022.csv")

# Related readers:
# readxl::read_excel("file.xlsx")
# object <- readRDS("object.rds")
\end{lstlisting}
\begin{itemize}
\item CSV: rectangular text data.
\item Excel: one or more worksheets.
\item RDS: one R object saved with its classes and attributes.
\end{itemize}
\end{frame}

\begin{frame}[fragile]{Inspect before transforming}
\begin{lstlisting}[language=R]
glimpse(health)
names(health)
count(health, income_level_current)
summary(select(health, gdp_per_capita_ppp,
                       under5_mortality))
n_distinct(health$iso3c)
\end{lstlisting}
The first reading should establish the row unit, the key, the variables, their types, and missing coverage.
\end{frame}

\begin{frame}[fragile]{Select columns and filter rows}
\begin{lstlisting}[language=R]
health_focus <- health |>
  select(country, region, income_level_current,
         gdp_per_capita_ppp, under5_mortality) |>
  filter(gdp_per_capita_ppp > 0)
\end{lstlisting}
\small \texttt{select()} changes columns; \texttt{filter()} changes rows. Neither computes a group statistic.
\end{frame}

\begin{frame}[fragile]{Construct and order variables}
\begin{lstlisting}[language=R]
health_focus <- health_focus |>
  mutate(log_income = log(gdp_per_capita_ppp)) |>
  arrange(desc(under5_mortality))
\end{lstlisting}
\texttt{mutate()} preserves the row unit while adding or changing columns. \texttt{arrange()} changes display order; it does not change the sample.
\end{frame}

\begin{frame}[fragile]{Summarise within groups}
\begin{lstlisting}[language=R]
health_table <- health_focus |>
  group_by(income_level_current) |>
  summarise(
    economies = n(),
    median_income = median(gdp_per_capita_ppp),
    median_mortality = median(under5_mortality),
    .groups = "drop"
  )
\end{lstlisting}
The input has one row per economy. The output has one row per income group.
\end{frame}

\begin{frame}[fragile]{Map variables to a figure}
\begin{lstlisting}[language=R]
ggplot(health_focus,
       aes(x = gdp_per_capita_ppp,
           y = under5_mortality,
           color = income_level_current)) +
  geom_point(alpha = 0.75) +
  scale_x_log10() +
  labs(x = "GDP per capita, PPP (log scale)",
       y = "Deaths per 1,000 live births",
       color = "Current income group") +
  theme_minimal()
\end{lstlisting}
\end{frame}

\begin{frame}{Checkpoint 1: choose the sequence}
You need the three economies in each income group with the highest under-five mortality in 2022.
\vspace{0.3cm}
\begin{enumerate}
\item What should one row of the final object represent?
\item Which functions would you use, and in what order?
\item What information would disappear if you summarised too early?
\end{enumerate}
\vfill
\centering\textbf{2 minutes discuss \(\longrightarrow\) 2 minutes share}
\end{frame}

\begin{frame}{Checkpoint 1: resolution}
One defensible sequence is:
\vspace{0.3cm}
\begin{center}
\texttt{filter(year == 2022)} \(\rightarrow\)
\texttt{group\_by(income group)} \(\rightarrow\)
\texttt{slice\_max(mortality, n = 3)}
\end{center}
\vspace{0.4cm}
The resulting row remains an economy. A group-level \texttt{summarise()} would remove economy names and prevent the requested ranking.
\end{frame}

\begin{frame}{Cross-section or endpoint comparison?}
\small
\begin{tabularx}{\textwidth}{p{2.6cm} X X}
\toprule
Design & Unit & Appropriate question \\
\midrule
2022 cross-section & economy in 2022 & How do economies differ at one date? \\
Fixed endpoints & economy with the same two years observed & How much did each economy's measure change between two dates? \\
Country-year panel & economy-year & How do levels and changes evolve over time? \\
\bottomrule
\end{tabularx}
\vspace{0.35cm}
The health and gender files are cross-sections. The access and climate files use fixed endpoints.
\end{frame}

\begin{frame}{Break}
\centering
\vfill
{\Huge\color{crimson}10 minutes}
\vspace{0.5cm}

When we return, we will place the health analysis in a Quarto document and render it from a clean session.
\vfill
\end{frame}
